Producers of electrical and electronic equipment in India must meet recycling obligations under the
E-Waste (Management) Rules, 2022 by buying Extended Producer Responsibility (EPR) credits, and the
compliance officer must show that the obligation, the certificates, the transfers, the payments and the
portal balance agree. The supporting records are spread across PDFs, spreadsheets and portal exports and
disagree in specific ways, so the officer must decide which discrepancy to investigate next and what
evidence supports a corrective action. This report presents \emph{ReCycleGuard}, a local, human-reviewed
assistant for that decision. A bounded LangGraph agent alternates between planning a check and assessing
its result; it may call six audited, read-only tools, and every arithmetic result comes from deterministic
exact-decimal code rather than from the model. Model-proposed findings are accepted only through a gate
that checks citations, numbers and prohibited claims, and every run ends in a report that a human
reviewer must decide. The system was built for a pinned 4-bit Qwen3.5-9B model served on the same machine,
and it was evaluated on a synthetic benchmark of 48 dossiers covering 16 fault families. On the 24
development dossiers, the final prompt completed 48 of 48 reliability runs without failure, with every
model reply valid against its schema and every sandbox probe denied. The development runs also show a
weakness: the model's single accepted finding in 48 runs was a false positive, and the fixed-retrieval variant averaged
134\,s per run against a 120\,s target. The held-out comparison of fixed retrieval with adaptive
investigation on 24 sealed dossiers is \pendingtext{held-out result}. The main limitations are the
synthetic data, a two-source regulatory corpus and a single local model.

\medskip
\noindent\textbf{Keywords:} agentic AI, extended producer responsibility, e-waste, retrieval-augmented
generation, plan-and-assess loop, human-in-the-loop, auditability.
